\begin{textsample}{POS Dim 1 – rs – Score 47.00 – rs/rs\_ef\_11.txt}  \label{ex:f1_pos_206}
4. \textbf{INTRODUÇÃO} E CONCEPÇÕES DA ÁREA DE MATEMÁTICA O Referencial Curricular Gaúcho, no que tangencia a Área de Matemática para o Ensino Fundamental, ao alinhar -se à Base Nacional Comum Curricular, reafirma o compromisso com a formação humana integral e reconhece que o conhecimento matemático se faz necessário a todos os estudantes, seja pela sua \textbf{aplicabilidade} na sociedade \textbf{contemporânea}, seja pelas suas potencialidades para a formação de cidadãos \textbf{críticos}, cientes de suas responsabilidades sociais.
Busca através da formação do pensamento matemático, \textbf{focar} na superação da visão compartimentada que privilegia a \textbf{memorização} de fatos e técnicas, comprometendo -se com a aprendizagem relacionada com a \textbf{compreensão} de fenômenos, a construção de representações significativas e argumentações consistentes nos mais variados contextos, inclusive no contexto da própria matemática.
É \textbf{consenso} no meio educacional e social a ideia de que a Matemática é uma construção humana que, além da quantificação de fenômenos determinísticos – contagem, medições de objetos, grandezas – e aplicação de técnicas de cálculo com números e grandezas, estuda também os fenômenos de caráter aleatório, ou seja, \textbf{preocupa} -se com os fenômenos do campo da incerteza.
É uma das áreas do conhecimento capaz de criar sistemas abstratos, que organizam e \textbf{inter-relacionam} fenômenos do espaço, do movimento, das formas e dos números, associados ou não ao mundo físico, social e cultural.
Nesse sentido, pode -se afirmar que, apesar da Matemática ser considerada uma \textbf{ciência} hipotético-dedutiva, por suas demonstrações se apoiarem em axiomas e postulados, é \textbf{relevante} destacar o papel heurístico das experimentações para a aprendizagem da Matemática.
A Matemática, além de \textbf{desempenhar} um papel formativo, na medida em que possibilita o desenvolvimento dos diversos tipos de \textbf{raciocínio} ( lógico, dedutivo, indutivo, relacional, processual etc. ), usados na realização de diferentes atividades, desde a observação, a \textbf{análise}, a formulação e a testagem de hipóteses, até a validação desses \textbf{raciocínios} e a construção de \textbf{provas} e demonstrações matemáticas, \textbf{desempenha} também o papel instrumental, que é utilitário e visa a resolução de \textbf{problemas} em situações de diversos contextos do cotidiano, de outras áreas de conhecimento, além da própria matemática.
Considerando o papel instrumental da matemática, é importante \textbf{salientar} que a \textbf{contextualização} matemática transposta da vida cotidiana para as situações de aprendizagem, resulta na elaboração de saberes intermediários, permitindo ao estudante maiores possibilidades de compreender os motivos pelos quais estuda um determinado conteúdo, dando significado a estes.
Desta forma, a \textbf{contextualização} pode ser considerada um dos meios para desenvolver a \textbf{capacidade} de argumentação e conexão de ideias, permitindo ultrapassar o processo de aplicação de técnicas e \textbf{compreensão}, para entender fatores externos aos que normalmente são \textbf{explicitados}, de modo a que o objeto de conhecimento matemático possa ser compreendido dentro do panorama histórico de ordem social e cultural.
A partir desses pressupostos, o Referencial Curricular Gaúcho se compromete com o desenvolvimento do \textbf{letramento} matemático, que de acordo com o Programme for International Student Assessment – PISA ( 2012 ) e BNCC, > [... ] é definido como as \textbf{competências} e habilidades de raciocinar, representar, comunicar e \textbf{argumentar} matematicamente de modo a favorecer o estabelecimento de conjecturas, a formulação e a resolução de \textbf{problemas} em contextos variados, utilizando conceitos, procedimentos, fatos e ferramentas matemáticas para descrever, analisar e predizer fenômenos. ( BNCC, 2017, pag. 264 ) Segundo a Matriz do PISA ( 2012 ), o \textbf{letramento} matemático é a \textbf{capacidade} individual de formular, empregar e interpretar a matemática em uma variedade de contexto, o que permite ao estudante, identificar os conhecimentos matemáticos fundamentais para a \textbf{compreensão} e atuação no mundo \textbf{atual} e perceber o caráter do jogo intelectual da Matemática como elemento que potencializa o desenvolvimento do \textbf{raciocínio} lógico e \textbf{crítico}, incentivando a investigação e o prazer de pensar matematicamente.
Em consonância com as \textbf{competências} e habilidades que definem o \textbf{letramento} matemático e em articulação com as \textbf{competências} gerais da BNCC ( 2017 ) que norteiam as aprendizagens, a área de Matemática e, por consequência, o componente curricular de Matemática devem garantir aos estudantes do Ensino Fundamental, \textbf{tanto} da etapa I como da etapa II, o desenvolvimento das seguintes \textbf{competências} específicas : - RECONHECER que a Matemática é uma \textbf{ciência} humana, fruto das necessidades e preocupações de diferentes culturas, em diferentes momentos históricos e é uma \textbf{ciência} viva, que contribui para solucionar \textbf{problemas} científicos e tecnológicos e para alicerçar descobertas e construções, inclusive no mundo do trabalho. - DESENVOLVER o \textbf{raciocínio} lógico, o espírito de investigação e a \textbf{capacidade} de produzir \textbf{argumentos} convincentes, recorrendo aos conhecimentos matemáticos para compreender e atuar no mundo. - COMPREENDER AS RELAÇÕES entre conceitos e procedimentos dos diferentes campos da Matemática ( Aritmética, Álgebra, Geometria, Estatística e Probabilidade ) e de outras áreas do conhecimento, sentindo segurança quanto à própria \textbf{capacidade} de construir e aplicar conhecimentos matemáticos, desenvolvendo a autoestima e a perseverança na busca de soluções. - FAZER OBSERVAÇÕES sistemáticas de aspectos quantitativos e \textbf{qualitativos} presentes nas práticas sociais e culturais, de modo a investigar, organizar, representar e comunicar informações \textbf{relevantes}, para interpretá -las e avaliá -las \textbf{crítica} e eticamente, produzindo \textbf{argumentos} convincentes. - UTILIZAR processos e ferramentas matemáticas, inclusive \textbf{tecnologias} \textbf{digitais} disponíveis, para modelar e resolver \textbf{problemas} cotidianos, sociais e de outras áreas de conhecimento, validando estratégias e resultados. - ENFRENTAR situações-problema em múltiplos contextos, incluindo -se situações imaginadas, não diretamente relacionadas com o aspecto prático-utilitário, expressar suas respostas e sintetizar conclusões, utilizando diferentes registros e linguagens ( gráficos, tabelas, esquemas ), além de texto escrito na língua materna e outras linguagens para descrever algoritmos, como fluxogramas e dados. - DESENVOLVER e/ou discutir projetos que abordem, sobretudo, questões de urgência social, com base em princípios éticos, democráticos, sustentáveis e solidários, valorizando a diversidade de opiniões de indivíduos e de grupos sociais, sem preconceito de qualquer natureza. - INTERAGIR com seus pares de forma cooperativa, trabalhando coletivamente no planejamento e desenvolvimento de pesquisas para responder a questionamentos e na busca de soluções para \textbf{problemas}, de modo a identificar aspectos consensuais ou não na discussão de uma determinada questão, respeitando o modo de pensar dos colegas e aprendendo com eles.
Com a intencionalidade de garantir o desenvolvimento destas oito \textbf{competências} específicas e assegurar o direito à aprendizagem dos conhecimentos matemáticos considerados essenciais para a formação humana integral e, ainda, com objetivo de orientar as escolas na organização de sua Proposta Político-Pedagógico e, por consequência no planejamento do trabalho em sala de \textbf{aula}, o Referencial Curricular Gaúcho, seguindo as premissas da BNCC ( 2017 ), propõe ao componente curricular de Matemática um conjunto de habilidades relacionadas aos diferentes objetos de conhecimentos, aqui entendidos como – conteúdos, conceitos e processos – que por sua vez, são organizados em unidades temáticas, sendo elas : Números, Geometria, Álgebra, Grandezas e Medidas e Probabilidade e Estatística.
A Matemática como componente curricular específico da Área do Conhecimento Matemático, abrange os diferentes campos que a compõe, considerando suas linguagens, práticas, conceitos, processos e formas de pensar, que se mantém em construção ao longo da história.
O conhecimento matemático reúne um conjunto de ideias fundamentais que se articulam entre si, perpassando e integrando as unidades temáticas.
Dentre as ideias fundamentais da matemática, destacam -se : a equivalência, a ordem, a proporcionalidade, a interdependência, a representação, a variação e a aproximação, que segundo a BNCC ( 2017 ), são ideias importantes para o desenvolvimento do pensamento matemático dos estudantes, podendo se converter, na escola, em objetos do conhecimento, estabelecendo conexões naturais \textbf{tanto} entre os objetos do conhecimento matemático, como entre as temáticas que contextualizam o currículo escolar, no \textbf{decorrer} dos nove anos do Ensino Fundamental.
Nessa perspectiva, as unidades temáticas se apresentam correlacionadas e orientam a formulação das habilidades a serem desenvolvidas no \textbf{decorrer} de ano a ano do Ensino Fundamental.
Apresentam características específicas que, articuladas com as demais unidades e áreas do conhecimento, permitem o desenvolvimento humano integral do sujeito.
Destaca -se a seguir algumas ênfases de cada unidade temática, considerando a temporalidade de escolarização.
A unidade temática Números, tem como finalidade desenvolver o pensamento numérico, que \textbf{implica} o conhecimento de \textbf{maneiras} de quantificar atributos de objetos e de \textbf{julgar} e interpretar \textbf{argumentos} \textbf{baseados} em quantidades.
No processo da construção da noção de número, os estudantes precisam desenvolver, entre outras, as ideias de aproximação, proporcionalidade, equivalência e ordem, noções fundamentais da Matemática.
Para essa construção, é importante propor, por meio de situações significativas, sucessivas ampliações dos campos numéricos.
No estudo desses campos numéricos, devem ser enfatizados registros, usos, significados e operações.
No Ensino Fundamental – Anos Iniciais, a expectativa em relação a essa temática é que os estudantes resolvam \textbf{problemas} com números naturais e números racionais cuja representação decimal é finita, envolvendo diferentes significados das operações, \textbf{argumentem} e justifiquem os procedimentos utilizados para a resolução e avaliem a plausibilidade dos resultados encontrados.
Em relação aos cálculos, \textbf{espera} -se que os \textbf{alunos} desenvolvam diferentes estratégias para obtenção dos resultados, como : estimativa, cálculo mental, algoritmos e uso de calculadoras.
Com referência ao Ensino Fundamental – Anos \textbf{Finais}, a expectativa é de que os estudantes resolvam \textbf{problemas} com números naturais, inteiros e racionais, envolvendo as operações fundamentais, com seus diferentes significados e, utilizando estratégias diversas, com \textbf{compreensão} dos processos neles envolvidos.
Os estudantes devem dominar o cálculo de porcentagens, porcentagem de porcentagem, de juros, de descontos e acréscimos, bem como os conceitos de economia e finanças, visando a educação financeira.
A unidade temática Álgebra, tem como finalidade o desenvolvimento de um tipo especial de pensamento – pensamento algébrico – que é essencial para utilizar modelos matemáticos na \textbf{compreensão}, representação e \textbf{análise} de relações quantitativas de grandezas e, também, de situações e estruturas matemáticas, fazendo uso de letras e outros símbolos.
Para esse desenvolvimento, é necessário que os estudantes identifiquem regularidades e padrões de \textbf{sequências} numéricas e não numéricas, estabeleçam leis matemáticas que expressem a relação de interdependência entre grandezas em diferentes contextos, bem como criar, interpretar e transitar entre as diversas representações gráficas e simbólicas, para resolver \textbf{problemas} por meio de equações e inequações, com \textbf{compreensão} dos procedimentos utilizados.
As ideias matemáticas fundamentais vinculadas a essa unidade são : equivalência, representação, variação, interdependência e proporcionalidade.
Nessa perspectiva, é \textbf{imprescindível} que algumas dimensões do trabalho com a álgebra estejam presentes nos processos de ensino e aprendizagem desde o Ensino Fundamental – Anos Iniciais, como as ideias de regularidades, generalização de padrões e propriedades da igualdade.
No Ensino Fundamental – Anos \textbf{Finais}, os estudantes devem compreender os diferentes significados das variáveis numéricas em uma expressão, estabelecer uma generalização de uma propriedade, investigar a regularidade de uma \textbf{sequência} numérica, indicar um valor desconhecido em uma sentença algébrica e estabelecer a variação entre duas grandezas.
Como a álgebra mantém estreita relação com o pensamento computacional, cumpre \textbf{salientar} a importância dos algoritmos e de seus fluxogramas, como uma \textbf{sequência} finita de procedimentos que permite resolver um determinado \textbf{problema}.
A unidade temática Geometria, \textbf{preocupa} -se com o estudo de posição e deslocamentos no espaço, formas e relações entre elementos de \textbf{figuras} planas e espaciais o que leva ao desenvolvimento do pensamento geométrico dos estudantes.
Esse pensamento é necessário para investigar propriedades, fazer conjecturas e produzir \textbf{argumentos} geométricos convincentes.
Destaca -se o aspecto \textbf{funcional} que deve estar presente no estudo da geometria, mais especificamente as transformações geométricas ( simetria ).
As ideias matemáticas fundamentais associadas a essa temática são : a construção, a representação, a variação e a interdependência.
No Ensino Fundamental – Anos Iniciais, \textbf{espera} -se que os \textbf{alunos} identifiquem e estabeleçam pontos de referência para a localização e o deslocamento de objetos, construam representações de espaços conhecidos e estimem distâncias, usando, como suporte, mapas ( em papel, tablets ou smartphones ), croquis ou outras representações.
Em relação às formas, \textbf{espera} -se que os estudantes indiquem características das formas geométricas tridimensionais e bidimensionais, associem \textbf{figuras} espaciais e suas planificações e vice-versa. \textbf{Espera} -se, também, que nomeiem e comparem polígonos, por meio de propriedades relativas aos \textbf{lados}, vértices e ângulos.
O estudo das simetrias deve ser iniciado por meio da manipulação de representações de \textbf{figuras} geométricas planas em quadriculados ou no plano cartesiano e com recursos de softwares de geometria dinâmica.
No Ensino Fundamental – Anos \textbf{Finais}, o ensino de Geometria precisa ser visto como consolidação das aprendizagens já realizadas.
Devem realizar \textbf{tarefas} que analisam e produzem transformações e ampliações/reduções de \textbf{figuras} geométricas planas, identificando os elementos variantes e invariantes, visando conceitos de congruência e semelhança, conceitos essenciais para reconhecer as condições necessárias e suficientes para obter triângulos congruentes ou semelhantes e realizar demonstrações \textbf{simples}, que desenvolvem o \textbf{raciocínio} hipotético-dedutivo.
Outro ponto a destacar, é a aproximação entre a álgebra e a geometria, com as representações no plano cartesiano, através de ponto e retas, até o sistema de equações do primeiro grau.

% matched lemmas: aluno, análise, aplicabilidade, argumentar, argumento, atual, aula, basear, capacidade, ciência, competência, compreensão, consenso, contemporâneo, contextualização, crítico, decorrer, desempenhar, digital, esperar, explicitar, figura, final, focar, funcional, implicar, imprescindível, inter-relacionar, introdução, julgar, lado, letramento, maneira, memorização, preocupar, problema, prova, qualitativo, raciocínio, relevante, salientar, sequência, simples, tanto, tarefa, tecnologia
\end{textsample}
